\documentclass[letterpaper,11pt]{article}
\usepackage[top=0.8in, bottom=0.8in, left=0.9in, right=0.9in]{geometry}
\usepackage[hidelinks]{hyperref}
\usepackage{lmodern}
\usepackage{parskip}
\usepackage{microtype}
\setlength{\parskip}{6pt}
\setlength{\parindent}{0pt}
\begin{document}
\begin{flushleft}
\textbf{\large Ajay Venkatesh} \\
Brooklyn, NY \\
+1 516 585 9013 \\
\href{mailto:av3855@nyu.edu}{av3855@nyu.edu}
\end{flushleft}
\vspace{10pt}
May 21, 2026
\vspace{6pt}

Hiring Manager \\
Dana-Farber Cancer Institute

\vspace{10pt}

Dear Hiring Manager,

My name is Ajay Venkatesh. I'm finishing my M.S. in Computer Engineering at NYU, with production software experience at PwC in Bengaluru, a summer SDE internship at Amazon, and ongoing on-campus work at NYU's Novel AI Technologies. I'm writing about the entry-level Bioinformatics Engineer role supporting thoracic cancer research at Dana-Farber.

The role's framing of designing software, databases, and interfaces for biological data, plus contributing to AI-driven platforms for clinical research, sits where my engineering and applied AI work have been pushing. At Amazon I deployed a \textbf{Retrieval-Augmented Generation (RAG)} pipeline with an \textbf{agentic workflow} that reasons over operational data, and implemented a \textbf{Model Context Protocol (MCP)} server for tool-augmented retrieval. I've also built RAG systems on \textbf{AWS Bedrock} with FAISS indexing, fine-tuned RoBERTa with custom LoRA, and shipped \textbf{ETL pipelines} at NYU transforming \textbf{NoSQL} into \textbf{PostgreSQL} with concurrency control for healthcare data.

On healthcare and clinical research collaboration: at NYU's Novel AI Technologies I architected an \textbf{NSF-funded} AI-powered health platform for iOS and Android, ingesting real-time wearable telemetry and deploying cloud-hosted AI models for health-metric analysis. At PwC I owned a distributed backend on \textbf{Microsoft Azure} for a healthcare platform, integrating with \textbf{Microsoft CRM}, with event-driven processing and fault-tolerant retries over \textbf{SQL Server}. Working alongside clinical and research teams in domains where the engineering decisions translate into patient outcomes is exactly where I've spent time.

On data-mining and software for biological systems: my graduate work includes \textbf{Machine Learning}, \textbf{Deep Learning}, and \textbf{Big Data} coursework, plus computer-vision pipelines through YOLOv8 and EfficientNet for image-based health workflows. I'm comfortable with the kind of mixed structured / unstructured data analysis bioinformatics requires.

Stack-wise: \textbf{Python}, \textbf{Java}, \textbf{SQL}, \textbf{Pandas}, \textbf{NumPy}, \textbf{PyTorch}, \textbf{Hugging Face}, plus \textbf{AWS} (SageMaker, Bedrock, ECS), \textbf{Docker}, and \textbf{PostgreSQL / MongoDB} for databases. Comfortable across the production-software and applied-ML boundary that bioinformatics work spans.

What pulls me toward Dana-Farber is the seriousness of the mission. Cancer research engineering runs on accuracy and reliability that real patients depend on, where careful software, careful data handling, and collaboration with clinical researchers compound into meaningful outcomes.

I'd welcome the chance to discuss further. Thanks for considering my application.

\vspace{12pt}
Sincerely, \\
Ajay Venkatesh

\end{document}